\chapter{Two More Worlds}\label{ch:worlds}

The engine's term type, its rewriter, its printer, its origin tracker and its notebook were built
for algebra. The last week of the project asked whether they were built for anything else, by
adding two small worlds from the author's other books: the untyped $\lambda$-calculus, and finite
partial orders with joins, meets and fixed points. The answer is yes, with one budget and a
different shape of theorem.

\section{The lambda calculus}

A cell is a $\lambda$-cell if it contains a $\lambda$ or a backslash, a \texttt{:=}, or starts with
a defined name; the engine decides. The parser reads $\lambda x\,y.\,e$, juxtaposition, digits as
Church numerals, and \texttt{name := term}; the Church library --- \texttt{succ}, \texttt{add},
\texttt{mul}, \texttt{true}, \texttt{false}, \texttt{if}, pairs --- is preloaded.

Reduction is normal order, one $\beta$-step at a time, each step recorded:

\begin{minted}{lean}
def betaStep : Term → Option (Term × Bool)
  | .app (.lam x body) arg => some (subst x arg body)
  | .app a b =>
    match betaStep a with
    | some (a', r) => some (.app a' b, r)
    | none => (betaStep b).map fun (b', r) => (.app a b', r)
  | .lam x e => (betaStep e).map fun (e', r) => (.lam x e', r)
  | .var _ => none
\end{minted}

The Boolean says whether an $\alpha$-renaming preceded the step, so that the notebook can label
it \rn{lambda.alpha-beta} rather than \rn{lambda.beta}.

\paragraph{Substitution had to be total.} The textbook capture-avoiding substitution renames a
binder and recurses on the \emph{renamed} body, and Lean cannot see that terminate --- the
author's earlier book marks it \lc{partial}. The engine carries the renaming down in one
structural pass first (\lc{freshen}, with the clash set and the renaming as parameters) and then
substitutes without ever renaming. It computes the same thing, and it is total.

\paragraph{The one budget.} Whether a term has a normal form is undecidable, so reduction runs on
fuel: $\Omega$ is refused after 1000 $\beta$-steps with what it had become. It is the only budget
in the engine, and it is the honest one; everything in Part~III was about removing the others.

\paragraph{Encoded, not extended.} Terms are encoded into \lc{Expr} --- \lc{fn "λ" [var x, body]}
and \lc{fn "@" [f, a]} --- so selection, explanation and origin tracking work unchanged and the
printer learned two heads. The de Bruijn view is computed alongside every step and sent with it;
the notebook's View menu toggles a rendering the engine already delivered. A normal form that is a
Church numeral or a Boolean is read out beside the result, and the reader is proved right on the
engine's own numerals:

\begin{minted}{lean}
theorem readChurch_church (n : Nat) : readChurch (church n) = some n
\end{minted}

What is not proved --- that the renaming preserves $\alpha$-equivalence, that the de Bruijn view
commutes with $\beta$ --- is why the $\beta$-steps are reported \status{unverified}. The
substitution lemma that \emph{is} proved says only that substitution introduces no new free
variables. Appendix~\ref{app:rules} lists all of it.

\section{Finite partial orders}

The order world is a different shape of engine: decisions over lists, and theorems that say each
decision means the textbook Prop. A poset is given as elements and a relation,
\texttt{poset(\{a,b,c\}; a<b, a<c)}; the engine takes the reflexive-transitive closure and decides
reflexivity, antisymmetry and transitivity, naming a witness on failure. Then: covers (the Hasse
diagram), upper and lower bounds, join and meet, lattice-or-witness-pair, top and bottom, maximal
and minimal, \texttt{le} explained by a chain of covers; maps as tables, \texttt{monotone} with a
witness; \texttt{lfp} and \texttt{gfp} by iterating from $\bot$ or $\top$, with the Kleene chain as
the derivation. \texttt{divisors(n)}, \texttt{subsets(\{…\})} and \texttt{chain(n)} build the
standard examples.

The theorems, all Init-only:

\begin{minted}{lean}
theorem checkPartialOrder_none {P : Poset} (h : checkPartialOrder P = none) :
    (∀ x ∈ P.elems, P.rel x x = true) ∧
    (∀ p ∈ P.le, p.1 ≠ p.2 → P.rel p.2 p.1 = false) ∧
    (∀ p ∈ P.le, ∀ q ∈ P.le, p.2 = q.1 → P.rel p.1 q.2 = true)

theorem covers_spec (P : Poset) (x y : String) :
    covers P x y = true ↔ P.lt x y = true ∧ ∀ z ∈ P.elems, ¬ (P.lt x z = true ∧ P.lt z y = true)

theorem sup_spec {P : Poset} {xs : List String} {j : String} (h : sup P xs = some j) :
    j ∈ upperBounds P xs ∧ ∀ v ∈ upperBounds P xs, P.rel j v = true

theorem iter_le_fixed {P : Poset} {f : PMap} (hm : monotoneFailure P f = none) {y : String}
    (hy : f.apply y = y) : ∀ (n : Nat) (x : String), P.rel x y = true → P.rel (f.iter n x) y = true
\end{minted}

The last is the Knaster--Tarski fragment a finite poset needs: every point of the Kleene chain
lies below every fixed point above its start, so with $x = \bot$ the chain's last element --- a
fixed point, as checked --- is the least. The honest boundary is that the chain stabilizes within
$|P|$ steps; that is pigeonhole on a finite chain, and it is checked at run time rather than
proved.

\begin{logbox}[The same shape as \lcn{checked}]
  Every theorem in this world is ``the decision means the Prop''. It is the shape of the
  run-time boundaries of Chapter~\ref{ch:order}, but here the Prop is the definition a reader
  meets in a first course, which is what should sit next to a Hasse diagram. The notebook draws
  the diagram in layers by height from a list of covers, and \lc{covers_spec} is the reason the
  picture is right.
\end{logbox}

\begin{trybox}
\begin{minted}{text}
let D = divisors(12)
join(D, 4, 6)
lattice(D)
let f = map(D; 1->2, 2->2, 3->6, 4->4, 6->6, 12->12)
lfp(D, f)
add 2 3
(λx. x x) (λx. x x)
\end{minted}
  The Kleene chain is the derivation of \texttt{lfp}. Type \texttt{\textbackslash lam} and
  Tab for the $\lambda$. The result of \texttt{add 2 3} reads ``the Church numeral 5''; toggle
  de Bruijn indices in the View menu. The last line is $\Omega$, refused after 1000 steps with
  what it had become.
\end{trybox}

\section{Exercises}

\begin{exercisebox}[The renaming pass]
  Show that \lc{freshen} followed by \lc{substRaw} computes the same term as the textbook's
  recursive-renaming substitution on $(\lambda y.\, x\, y)[x := y]$, and say what the textbook
  version recurses on that Lean cannot see decrease.
\end{exercisebox}

\begin{exercisebox}[Pigeonhole, proved]
  State the theorem that closes the order world's boundary: on a finite poset, a monotone map
  iterated from $\bot$ reaches a fixed point within $|P|$ steps. Which lemma about
  \lc{List.Nodup} would you reach for?
\end{exercisebox}
